\usepackage[a4paper,top=2.7cm,bottom=2.7cm,inner=2.8cm,outer=2.4cm]{geometry}
\usepackage{fontspec}
\usepackage{xeCJK}
\usepackage{microtype}
\usepackage{amsmath,amssymb,mathtools}
\usepackage{booktabs,tabularx,longtable}
\usepackage{xcolor}
\usepackage{enumitem}
\usepackage{csquotes}
\usepackage{hyperref}
\usepackage{bookmark}
\usepackage{tcolorbox}
\tcbuselibrary{breakable,skins}
\usepackage{makeidx}
\usepackage[backend=biber,style=authoryear,sorting=nyt,maxcitenames=2,maxbibnames=99,
  uniquename=false,giveninits=true,date=year]{biblatex}

\makeindex
\hypersetup{colorlinks=true,linkcolor=black,citecolor=blue!45!black,urlcolor=blue!55!black,
  pdftitle={我们如何知道？——问题驱动的认识论},
  % TODO（出版前）：pdfauthor 与 main.tex 的 \author 一并补齐。
  pdfauthor={}}
\setlength{\parindent}{2em}
\setlength{\parskip}{0.22em}
\emergencystretch=2em
\setlist{nosep,leftmargin=2.2em}
\renewcommand{\arraystretch}{1.25}
\definecolor{caseblue}{RGB}{236,244,250}
\definecolor{arggray}{RGB}{245,245,242}
\definecolor{comparegreen}{RGB}{238,247,240}
\definecolor{exerciseorange}{RGB}{253,246,234}

\newtcolorbox{openingcase}{breakable,colback=caseblue,colframe=blue!35!black,
  title={开篇案例},fonttitle=\bfseries}
\newtcolorbox{argumentbox}[1][]{breakable,colback=arggray,colframe=black!45,
  title={论证重构 #1},fonttitle=\bfseries}
\newtcolorbox{comparison}{breakable,colback=comparegreen,colframe=green!35!black,
  title={比较视角},fonttitle=\bfseries}
\newtcolorbox{checkpoint}{breakable,colback=white,colframe=black!35,
  title={自测小结},fonttitle=\bfseries}
\newtcolorbox{exercises}{breakable,colback=exerciseorange,colframe=orange!55!black,
  title={章末练习},fonttitle=\bfseries}
\newcommand{\term}[2]{\textbf{#1}（\textit{#2}）\index{#1}}
% 段内问题提示：保留论证路标，但不再把每个细分点切成独立小节。
\newcommand{\minitopic}[1]{\par\medskip\noindent\textbf{#1}\enspace\ignorespaces}
\newcommand{\keyreading}[1]{\par\smallskip\noindent\textbf{核心阅读：}#1}
\newcommand{\furtherreading}[1]{\par\noindent\textbf{延伸阅读：}#1}
\newtcolorbox{readingmap}{breakable,colback=white,colframe=blue!28!black,
  title={分层阅读地图},fonttitle=\bfseries,before skip=1em,after skip=.5em}
\newcommand{\readingstrand}[3]{%
  \par\noindent\textbf{#1：}#2\space\parencite{#3}\par\smallskip}
\newcommand{\chaptergoals}[1]{\begin{tcolorbox}[breakable,colback=white,colframe=blue!30!black,
  title={学习目标},fonttitle=\bfseries]#1\end{tcolorbox}}

\ctexset{
  part={format=\centering\Huge\bfseries,name={第,部},number=\chinese{part}},
  chapter={format=\raggedright\Huge\bfseries,name={第,章},number=\chinese{chapter}},
  section={format=\Large\bfseries},
  subsection={format=\large\bfseries}
}
